% Target venue: IEEE CiFer 2027 -- Computational Intelligence for Financial
% Engineering and Economics, within IEEE SSCI 2027.
\documentclass[10pt,conference]{IEEEtran}
\usepackage{booktabs}
\usepackage{url}
\usepackage[hidelinks]{hyperref}
\usepackage{amsmath}
\usepackage{amssymb}
\usepackage{graphicx}

\title{Separating Language-Model Explanation from Portfolio Allocation:\\A Prospective Protocol and Three Retrospective Null Results}
\author{Author information withheld for double-blind review}

\begin{document}
\maketitle

\begin{abstract}
Language models are increasingly proposed as portfolio allocators, although historical prompts may reveal outcomes already embedded in model training. We present an auditable architecture in which a dated deterministic rule selects and weights assets, a language model only explains approved facts, and a human records approval. A separate retrospective stress test relaxes that boundary in four arms over thirteen years: named issuers, anonymised issuers, a deterministic factor control, and direct model-generated weights. The measured differences are indistinguishable from zero: named versus anonymised $+0.38$ pp/yr ($p=0.91$), anonymised versus deterministic $-0.05$ pp/yr ($p=0.99$), and solver-routed versus direct weights $+0.81$ pp/yr ($p=0.78$). Anonymisation cannot prove temporal purity because financial attributes may still identify firms and pretrained memory is unobservable. The direct-weight arm also failed to sum to one in five of thirteen years and omitted eligible names in two. The evaluation infrastructure contains 514 issuers, including 166 that stop trading, receipt-date gating of regulatory filings, a nested selection protocol with a 0.65 pp/yr hindsight premium, and a deflated-Sharpe probability of 0.986 over 36 trials. Seven audited defects had inflated earlier results. Accordingly, the historical evidence diagnoses the measurement system rather than validating language-model alpha; scientific evaluation of the language component begins with a frozen prospective protocol.
\end{abstract}

\begin{IEEEkeywords}
portfolio optimisation, large language models, backtest overfitting, deflated Sharpe ratio, survivorship bias, temporally available data, emerging markets
\end{IEEEkeywords}

\section{Introduction}

A language model asked for portfolio weights will produce them. Whether an accessible model is useful to an investment decision, and which responsibility it should receive, are separate questions. Historical evaluation is especially fragile: company names, accounting patterns or salient events may connect a prompt to outcomes embedded in training, even when the input table itself is dated correctly.

We therefore fix a production architecture in which the model cannot select an asset or set a weight. It receives a completed decision record and produces only explanation, risks and questions for human review. A separate retrospective diagnostic temporarily relaxes this contract to measure three comparisons that are usually conflated:

\begin{enumerate}
\item Contamination sensitivity: does the model do better when it can see company names than when it sees anonymised financial facts? The gap can reveal contamination but cannot establish its absence.
\item Added value: does the model beat a deterministic control that ranks the same universe by the same pre-declared factor? Without this arm, any result merely restates the factor.
\item Value of decoupling: does routing a bounded score through a convex solver beat letting the model emit weights directly? This is what the architecture is for.
\end{enumerate}

Our contribution is threefold. First, a typed responsibility boundary linking dated data, deterministic calculation, explanation and human approval. Second, three retrospective null results on Brazilian equities, explicitly denied the status of prospective language-model validation. Third, the evaluation infrastructure: panel construction, availability gating, multiple-testing correction and seven defects found in our own work. The production protocol was frozen on 16 August 2026; any predictive role for a model would require a new prospective arm. Upon acceptance, we will release the evaluation artefacts and versioned manifest in a permanent repository; during double-anonymous review, the same materials are supplied in an anonymised archive.

\section{Related work}

Markowitz's mean--variance formulation \cite{markowitz} remains the allocation substrate, with estimation error, concentration and turnover treated as first-class constraints rather than afterthoughts. Black and Litterman \cite{blacklitterman} established the pattern we follow structurally: qualitative views are combined with a mechanical allocator through a defined interface rather than replacing it. DeMiguel et al. \cite{demiguel} show that optimised portfolios frequently fail to dominate naive diversification out of sample, which motivates explicit caps and a frozen comparator. Novy-Marx and Velikov \cite{novymarx} show that frictions can erase anomaly profitability, which is why cost and tax are inside the evaluated protocol here rather than appended to it.

On inference, Bailey and L\'opez de Prado \cite{deflated} give the deflated Sharpe ratio, which corrects the observed statistic for the number of trials and the higher moments of the return distribution; Bailey et al. \cite{pbo} give the probability of backtest overfitting. Harvey et al. \cite{harvey} document how weak the standard significance bar is once the size of the search is acknowledged. These are not optional refinements for a paper that evaluates 36 configurations.

Recent evidence argues against treating backtested language-model performance as a single construct. Perlin et al. run 30,000 anonymised simulations over 1,522 firms and find no consistent advantage over 1/N or the S\&P 500 \cite{perlin}. FINSABER reports deterioration over broad universes, long horizons and changing regimes \cite{finsaber}. Pelster and Val use a live earnings experiment, illustrating why prospective observation is the stronger contamination control \cite{pelsterval}. Kim et al. find useful narrative analysis of anonymised statements \cite{kim}, a role that does not imply trading alpha. FINCON's hierarchy motivates explicit responsibility boundaries, not autonomous trading here \cite{fincon}. These results motivate both our deterministic control and the narrower production role assigned to language.

\section{Architecture}

\subsection{Production interface}

Let $x_t$ contain only observations admitted by decision date $t$. A deterministic selector $S$ produces eligibility, factor scores and rejection reasons; an allocator $A$ produces weights under the policy $P$:

\begin{equation}
 d_t=S(x_t),\qquad w_t=A(d_t,P),\qquad \sum_i w_{i,t}=1.
\end{equation}

Only after $w_t$ is sealed does the language model receive approved fields $(d_t,w_t)$ and return prose $e_t$. There is no edge from $e_t$ to $S$, $A$ or $P$. A human reviewer records acceptance or rejection, identity and rationale. This interface makes the central claim testable through the stored inputs, outputs and hashes, independent of prose quality.

The eleven-year strategy uses a hard top-$N$ cut on the factor score, one share class per issuer and score-proportional weights inside fixed ceilings. MVO remains an independently implemented comparator. The convex solver described below belongs only to the four-arm diagnostic and is not the published allocation rule.

\subsection{Pipeline}

Each January the system executes five stages in a fixed order. Stages one to four are deterministic, language is generated only after the portfolio is sealed, and the last stage is human.

\emph{(i) Dated universe.} The universe of record is the exchange's historical quotation archive \cite{b3cotahist}. \emph{(ii) Available-at-decision fundamentals.} Only regulatory filings received by the decision date may enter \cite{cvmitr,cvmdfp}. \emph{(iii) Screen.} Assets must clear liquidity, profitability, return on capital and solvency, with distinct quantities for banks. \emph{(iv) Ranking and allocation.} Survivors are ranked by declared factors, allocated within ceilings, and residual capital is held at the domestic overnight rate \cite{bcb}. \emph{(v) Explanation and review.} The sealed proposal, rejection reasons, weights, order sizes and estimated costs are explained, then accepted or rejected by a named person. No order is transmitted.

The economic hypothesis is fundamental-led rather than language-model-led. Quality and value are the primary accounting dimensions, twelve-month momentum is a market confirmation signal, and liquidity is an execution gate. Banks and operating companies use different accounting quantities. The model does not create this signal: in production architecture it converts approved facts into thesis, risk and review questions. MVO is an independent benchmark in the published record and an allocator only inside the four-arm experiment.

\section{Evaluation infrastructure}

A null result is a claim about an instrument as much as about a treatment. This section describes the instrument.

\subsection{Survivorship}

A public adjusted-price provider serves the companies that still exist, so a universe built from one silently deletes every company that was delisted, acquired or wound up. Rebuilding from the exchange archive over 2010--2025 yields 4{,}066 sessions and 514 issuers, of which 166 stop trading before December 2025; 77 of those disappear before 2020, and 58 had already fallen more than 60\% from their own peak when they stopped trading. These are precisely the worst returns in the sample.

Delisted tickers are not served by a conventional price provider. The initial round-ratio detector had 88.1\% precision but only 23.3\% recall against the available event archive. We therefore added a primary B3 event collector that archived 18,190 cash, share and manual events and covered 475 of 497 price-panel securities. This does not retroactively certify the published series: 22 codes remain unavailable through the current endpoint and 271 subscriptions or security conversions require explicit resolution. The performance panel remains research-grade until it is rebuilt from raw closes with complete event coverage.

\subsection{Availability gating}

Trailing-twelve-month quantities are bridged as annual plus current year-to-date minus comparative year-to-date, using the year-to-date column only. The interim filing also publishes the isolated quarter, and mixing the two silently corrupts every ratio from the second quarter onward. This is the kind of defect that produces a better backtest and no error message.

\subsection{Holding, cost and tax}

The book is bought in January and held. This is worth stating because the natural vectorised expression of a backtest, compounding the inner product of returns and weights, silently rebalances to fixed weights every session at zero cost, which is neither the protocol under evaluation nor anything an investor can execute. Costs are the published exchange fee plus a slippage term growing with the position's participation in average traded volume. Taxation follows the domestic regime, with gains taxed when a review actually realises them and the final year charged as a full liquidation.

\subsection{Choosing the rule without reading the answer}

Our first protocol failed this test, and the failure is instructive. A grid of candidate rules was evaluated, the best was published, and its rank moved from 57th on the declared training criterion to 1st on the holdout, which is only possible if the holdout was read.

The current protocol is nested. For decision year $t$, the 36 configurations are ranked using only years that closed before $t$, by Sharpe of the excess return over cash; year $t$ is then evaluated once. Changing configuration is charged a full-turnover rebalance. Ranking requires three closed years, so 2012--2014 is the selection window and 2015--2025 is the evaluation window, giving eleven decisions.

Three quantities are published with every result: the nested outcome, the full-sample winner, and the difference between them. That difference, the \emph{hindsight premium}, is what an author would gain by publishing the search winner as though it were attainable. On this panel it is 0.65 pp/yr. On a panel with one year less of training it was 4.98 pp/yr, which is a useful calibration of how much the measure moves with sample length.

With 36 trials, the winner's Sharpe is biased upward by the search itself. The observed excess-return Sharpe is 0.933 against an expected maximum under the null of 0.355, giving a deflated-Sharpe probability of 0.986, above the 95\% threshold.

\section{Experiment}

\subsection{Design}

Four retrospective arms use the same thirteen annual universes. The first three use the same constrained solver,

\begin{equation}
\max_w\;\mu(c)^\top w-\frac{\gamma}{2}w^\top\Sigma w-\kappa\lVert w-w_{\mathrm{prev}}\rVert_1,
\end{equation}

with $\sum_iw_i=1$, $w\ge0$ and fixed ceilings. The model score $c_i\in[-1,1]$ can only tilt $\mu_i$ inside this diagnostic. The production artifact never uses this channel.

\begin{itemize}
\item Named: the model sees company names and returns a bounded score.
\item Anonymised: identical, but issuers are identified only by integers.
\item Deterministic: control with no model, ranking by the pre-declared factor score.
\item Monolithic: the counterfactual in which the model returns weights directly, with no solver and no constraints.
\end{itemize}

Comparisons are paired by year and tested with a paired $t$-test.

\begin{table}[t]
\caption{Four arms, thirteen years, same eligible universe.}
\centering
\small
\begin{tabular}{lrrr}
\toprule
Arm & CAGR & Win yrs & Eq. wt. \\
\midrule
Named & 13.00\% & 7 of 13 & 43.8\% \\
Anonymised & 12.62\% & 10 of 13 & 44.7\% \\
Deterministic & 12.67\% & 8 of 13 & 52.2\% \\
Monolithic & 11.81\% & 7 of 13 & 55.0\% \\
\midrule
Cash (CDI) & 9.59\% & --- & --- \\
\bottomrule
\end{tabular}
\label{tab:arms}
\end{table}

\begin{table}[t]
\caption{The three effects. None survives testing.}
\centering
\begin{tabular}{lrr}
\toprule
Effect & Annualised & $p$ \\
\midrule
Contamination (named $-$ anonymised) & $+0.38$ pp & 0.91 \\
Added value (anonymised $-$ deterministic) & $-0.05$ pp & 0.99 \\
Decoupling (anonymised $-$ monolithic) & $+0.81$ pp & 0.78 \\
\bottomrule
\end{tabular}
\label{tab:effects}
\end{table}

\subsection{Results}

Table~\ref{tab:effects} reports three nulls, and each means something different.

\emph{No detectable named--anonymised gap.} This is not evidence of temporal purity. The model may infer identity from financial attributes, both arms may contain pretrained memory, and the training corpus is not observable. The comparison is a sensitivity test with low power, not a retrospective certificate against contamination.

\emph{No detectable added value.} The model reproduces the deterministic ranking rather than adding judgement. The gap is five hundredths of a percentage point, with a negative sign. This is an uncomfortable finding for the prevailing narrative about language models in investment management, which is a reason to publish it rather than a reason not to.

Keeping the model away from the weights did not cost performance. The monolithic arm breached no ceiling in thirteen years, a point in its favour that we report because the opposite would have been the convenient finding. What it did produce was arithmetically invalid: weights failing to sum to one in five of thirteen years, ranging from 0.92 to 1.017, and 65 and 83 eligible names silently omitted in 2021 and 2022. The argument for the solver is therefore not that it prevents a limit breach. It is that it returns an allocation requiring no repair, and that repair on a malformed weight vector is itself an unlogged decision.

Thirteen annual observations make these underpowered tests. The honest statement is that no effect was detected, not that none exists.

\subsection{Robustness and the separate product record}

Holding the rule and panel fixed, annual reselection compounded at 18.58\% before costs, versus 16.37\% quarterly and 13.92\% monthly. After tax, the paired gaps were 1.55 pp/yr ($p=0.38$) and 3.27 pp/yr ($p=0.31$); no faster cadence improved the selection. A retention buffer also raised carryover from 36\% to 42\% and reduced one-year holdings from 21 to 16, but its 0.33 pp/yr return difference was not significant ($p=0.14$). These are stability checks, not additional return claims.

The supplementary product record distinguishes the Ibovespa total-return index, which reinvests distributions in a theoretical portfolio \cite{b3ibov}, from BOVA11, an investable ETF subject to fees, trading costs and tracking difference \cite{blackrockbova}. A retrospective event-risk extension reduced maximum drawdown from 47.8\% to 28.7\%, but added no detectable return in 2019--2025 ($p=0.964$). It was designed after Covid-19 and omits intrayear tax, so it is not prospective validation and does not enter the language-model tests.

\section{Defects found by auditing our own backtest}

Seven defects were found and corrected. Every one inflated the result, which is what one should expect: an error that hurts your own numbers is noticed immediately and an error that helps them is not.

\begin{table}[t]
\caption{Defects, all of which had inflated the result.}
\centering
\footnotesize
\begin{tabular}{p{3.1cm}p{4.1cm}}
\toprule
Defect & Correction \\
\midrule
MVO comparator identical to the strategy & Independent implementation \\
Vendor panel dropped delisted issuers & Universe rebuilt from the exchange archive \\
Delisting year discarded entirely & Position liquidated into cash \\
No tax modelled & Domestic tax regime \\
Flat cost ignoring liquidity & Slippage by participation in volume \\
Implicit daily rebalancing & Exact buy-and-hold path \\
Session filter using a global peak & Local rolling peak \\
\bottomrule
\end{tabular}
\label{tab:defects}
\end{table}

After the final correction the published CAGR fell and the daily drawdown rose from 30.4\% to 47.8\%. The figures were published as they came out.

\section{Limitations}

The evaluation window is also the development window; the 0.65 pp hindsight premium bounds that problem without making the record prospective. Primary B3 events now cover 95.6\% of price-panel securities, but unresolved codes and security conversions still prevent institutional reconciliation, while issuers without provider coverage receive an explicitly listed cross-sectional median distribution yield. Eleven annual observations remain few: the deflated-Sharpe probability corrects search bias, not sample size. The model experiment covers one model and one training cutoff. The annual cadence beat quarterly and monthly reselection here, but is not universally optimal. Benevente~2 was designed after Covid-19 and omits intrayear tax, so it requires pre-registration and prospective evaluation. Timestamped news and event-driven decisions belong in a separate protocol. Factor transformations, thresholds, the configuration grid and annual choices are disclosed in the supplementary artefacts.

\section{Conclusion}

We measured a retrospective counterfactual for language-model discretion and detected no advantage: names did not help, the model did not beat its deterministic factor, and direct weights were malformed in five years. These nulls do not validate the model or prove contamination absent. They support the narrower architecture: dated data and deterministic code create the portfolio, language explains a sealed record, and a person remains responsible. Predictive claims require a prospective arm beginning after the frozen registration, with deterministic and anonymised controls retained.

\section*{Acknowledgment and generative-AI disclosure}
A generative-AI assistant was used for language editing and code review. The authors independently verified every calculation, citation and claim and accept full responsibility. Empirical results were generated by deterministic project code, not by the assistant \cite{openai}.

\clearpage
\begin{thebibliography}{00}
\bibitem{markowitz} H. Markowitz, ``Portfolio Selection,'' \emph{The Journal of Finance}, vol. 7, no. 1, pp. 77--91, 1952.
\bibitem{blacklitterman} F. Black and R. Litterman, ``Global Portfolio Optimization,'' \emph{Financial Analysts Journal}, vol. 48, no. 5, pp. 28--43, 1992.
\bibitem{demiguel} V. DeMiguel, L. Garlappi, and R. Uppal, ``Optimal versus Naive Diversification: How Inefficient is the 1/N Portfolio Strategy?'' \emph{The Review of Financial Studies}, vol. 22, no. 5, pp. 1915--1953, 2009.
\bibitem{novymarx} R. Novy-Marx and M. Velikov, ``A Taxonomy of Anomalies and Their Trading Costs,'' \emph{The Review of Financial Studies}, vol. 29, no. 1, pp. 104--147, 2016.
\bibitem{deflated} D. H. Bailey and M. L\'opez de Prado, ``The Deflated Sharpe Ratio: Correcting for Selection Bias, Backtest Overfitting, and Non-Normality,'' \emph{The Journal of Portfolio Management}, vol. 40, no. 5, pp. 94--107, 2014.
\bibitem{pbo} D. H. Bailey, J. Borwein, M. L\'opez de Prado, and Q. J. Zhu, ``The Probability of Backtest Overfitting,'' \emph{Journal of Computational Finance}, vol. 20, no. 4, pp. 39--69, 2016.
\bibitem{harvey} C. R. Harvey, Y. Liu, and H. Zhu, ``\ldots and the Cross-Section of Expected Returns,'' \emph{The Review of Financial Studies}, vol. 29, no. 1, pp. 5--68, 2016.
\bibitem{perlin} M. S. Perlin, C. R. Foguesatto, F. M. M\"uller, and M. B. Righi, ``Can AI Beat a Naive Portfolio? An Experiment with Anonymized Data,'' \emph{Finance Research Letters}, Art. no. 107126, 2025, doi: 10.1016/j.frl.2025.107126.
\bibitem{finsaber} Y. Li, K. Kim, M. Cucuringu, and S. Ma, ``Can LLM-Based Financial Investing Strategies Outperform the Market in Long Run?'' in \emph{Proc. 32nd ACM SIGKDD Conf.}, pp. 2711--2722, 2026, doi: 10.1145/3770854.3785702.
\bibitem{pelsterval} M. Pelster and J. Val, ``Can ChatGPT Assist in Picking Stocks?'' \emph{Finance Research Letters}, vol. 59, Art. no. 104786, 2024, doi: 10.1016/j.frl.2023.104786.
\bibitem{kim} A. G. H. Kim, M. Muhn, and V. V. Nikolaev, ``Financial Statement Analysis with Large Language Models,'' arXiv:2407.17866, 2024.
\bibitem{fincon} Y. Yu et al., ``FinCon: A Synthesized LLM Multi-Agent System with Conceptual Verbal Reinforcement for Enhanced Financial Decision Making,'' in \emph{Advances in Neural Information Processing Systems 37}, 2024, doi: 10.52202/079017-4354.
\bibitem{b3cotahist} B3 S.A. --- Brasil, Bolsa, Balc\~ao, ``Historical quotations,'' 2026, \href{https://www.b3.com.br/pt_br/market-data-e-indices/servicos-de-dados/market-data/historico/mercado-a-vista/cotacoes-historicas/}{online}.
\bibitem{b3ibov} B3 S.A. --- Brasil, Bolsa, Balc\~ao, ``Ibovespa B3,'' 2026, \href{https://www.b3.com.br/pt_br/market-data-e-indices/indices/indices-amplos/ibovespa.htm}{online}.
\bibitem{blackrockbova} BlackRock, ``iShares Ibovespa Fundo de {\'I}ndice (BOVA11): Factsheet,'' 2026, \href{https://www.blackrock.com/br/literature/fact-sheet/bova11-ishares-ibovespa-fundo-de-ndice-fund-fact-sheet-pt-lm.pdf}{online}.
\bibitem{cvmitr} Comiss\~ao de Valores Mobili\'arios, ``Formul\'ario de Informa\c{c}\~oes Trimestrais (ITR),'' CVM Open Data, \href{https://dados.cvm.gov.br/dados/CIA_ABERTA/DOC/ITR/DADOS/}{online}.
\bibitem{cvmdfp} Comiss\~ao de Valores Mobili\'arios, ``Demonstra\c{c}\~oes Financeiras Padronizadas (DFP),'' CVM Open Data, \href{https://dados.cvm.gov.br/dataset/cia_aberta-doc-dfp}{online}.
\bibitem{bcb} Banco Central do Brasil, ``Sistema Gerenciador de S\'eries Temporais, s\'erie 12 (CDI),'' \href{https://www3.bcb.gov.br/sgspub/}{online}.
\bibitem{openai} OpenAI, ``GPT-5 System Card,'' 2025, \href{https://openai.com/index/gpt-5-system-card/}{online}.
\end{thebibliography}
\end{document}
